# Theoretical Foundations of P-Scube

## 1. Problem Formulation: The Utility-Privacy Tension

In a graph stream $\mathcal{G}$, the degree distribution typically follows a power-law: $\Pr[d_v = k] \propto k^{-\gamma}$. 
Skew-aware summarization structures (like Scube) aim to minimize hash collisions by allocating resources $A_v$ (e.g., number of hash addresses) proportional to the estimated degree $\hat{d}_v$.

### Definition 1: Deterministic Resource Allocation (Original Scube)
In the original Scube, the resource allocation function is a deterministic step function:
$$ A_v = \max\left(2, \left\lceil \frac{\hat{d}_v}{\theta} \right\rceil \right) $$
where $\theta$ is the capacity threshold of a single row/column, and $\hat{d}_v = c \times E_{R,L}$ is the estimated degree from the probabilistic counter.

### Definition 2: Extractability Advantage
An adversary $\mathcal{A}$ observes the allocated resources $A_v$ (either directly or via timing side-channels) and attempts to infer if node $v$ is a "Hub" (i.e., true degree $d_v > \tau$). 
The **Extractability Advantage** is defined as:
$$ Adv(\mathcal{A}) = \left| \Pr[d_v > \tau \mid A_v] - \Pr[d_v > \tau] \right| $$
**Theorem 1 (Tension):** In a deterministic skew-aware summary, if $A_v$ strictly monotonically increases with $\hat{d}_v$, and the estimator $\hat{d}_v$ is unbiased, then for large $\tau$, $Adv(\mathcal{A}) \to 1 - \Pr[d_v > \tau]$. This means the adversary can perfectly extract hub nodes with near 100% confidence.

---

## 2. P-Scube Mechanism Design and Theoretical Guarantees

To mitigate this tension, P-Scube introduces two mechanisms that decouple the deterministic mapping while preserving the overall skew-aware utility.

### Mechanism 1: Monotonic Noisy Promotion
To prevent state oscillation in a streaming environment, we cannot add pure random noise at every step. Instead, we bind the noise to the node's unique identifier (fingerprint/hash).

Let $h(v)$ be the hash signature of node $v$. We generate a pseudo-random noise value $\eta_v$ drawn from a Laplace distribution:
$$ \eta_v \sim \text{Laplace}\left(0, \lambda \right) $$
*Implementation note: $\eta_v$ is deterministically generated using $h(v)$ as the seed, ensuring $\eta_v$ is constant for the same node across the entire stream.*

The noisy estimated degree is:
$$ \tilde{d}_v = \hat{d}_v + \eta_v $$

### Mechanism 2: Coarsened Resource Allocation
Instead of continuous expansion, we map the noisy degree to a set of discrete resource buckets $\mathcal{B} = \{B_1, B_2, \dots, B_k\}$ (e.g., $\{2, 4, 6\}$ addresses).
$$ A_v^* = \text{Bucketize}(\tilde{d}_v) $$
where $\text{Bucketize}(x) = B_i$ if $T_{i-1} \le x < T_i$.

### Theorem 2: Degree Indistinguishability (Local Differential Privacy)
By combining Monotonic Noisy Promotion and Coarsened Allocation, P-Scube satisfies $\epsilon$-Degree Indistinguishability. 
For any two nodes $u$ and $v$ whose true degrees differ by at most $\Delta$ (i.e., $|d_u - d_v| \le \Delta$), the probability that they are assigned to the same resource bucket $B_i$ satisfies:
$$ \frac{\Pr[A_u^* = B_i]}{\Pr[A_v^* = B_i]} \le e^{\epsilon} $$
where $\epsilon = \frac{\Delta}{\lambda}$.

**Proof Sketch:**
Since $\tilde{d} = \hat{d} + \eta$ and $\eta \sim \text{Laplace}(0, \lambda)$, the probability density function of $\tilde{d}$ satisfies the standard Differential Privacy property. The bucketization step (post-processing) does not increase the privacy budget $\epsilon$. Thus, the adversary cannot confidently distinguish between a node with degree $d$ and a node with degree $d + \Delta$, effectively hiding the exact rank of hub nodes and bounding $Adv(\mathcal{A})$.

---

## 3. Engineering Constraints for the Streaming Environment

1. **Monotonicity:** Because $\eta_v$ is fixed per node, as the true degree $d_v$ grows over time, the noisy degree $\tilde{d}_v$ grows monotonically. This guarantees that $A_v^*$ never decreases, which is a strict requirement for the Scube matrix (we cannot "shrink" a node's allocated addresses once edges are inserted).
2. **Low Overhead:** Generating $\eta_v$ via a fast hash-based pseudo-random number generator (PRNG) adds negligible latency compared to standard cryptographic DP mechanisms.